Un solide cubique, de masse $m=\qty{50,0}{\gram}$, est lancé avec
une vitesse de valeur $v=\qty{2,5}{\v}$.

Il glisse vers le haut sur une pente inclinée d'un angle $\alpha=\ang{15}$ sur
l'horizontale.

On néglige tous les frottements.

\begin{figure}[H]
    \centering
    \begin{tikzpicture}[solide/.style={minimum width=1cm,minimum height=6mm,draw,
            fill=black!20,thick}]
        \def\al{15}
        \def\L{6}
        \pgfmathparse{\L*tan(\al)}\let\H=\pgfmathresult
        \node[inner sep=0pt,minimum width=\L cm,minimum height=\H cm] (R) {};
        \draw (R.north east) coordinate(B) |- (R.south west) coordinate(A) -- cycle;
        \foreach \p/\v [count=\i] in {0/$\vec{v}$,0.7/$\vec{v}^{\prime}$}{
            \node[solide,above right=0.6pt and 0,
                  rotate=\al] (S\i) at ($(A)!\p!(B)$) {};
            \node[circle,inner sep=1pt,fill] at (S\i.center) {};
            \draw[->] (S\i.center) -- ++(\al:1.5) node[above] {\v};
        }
        \draw ([xshift=1.5cm]R.south west) arc (0:\al:1.5)
            node[right,pos=0.7] {$\alpha=\ang{\al}$};
        \draw[dashed] (S1.center) -- (S1.center -| S2.center) coordinate(tmp);
        \draw[<->,shorten <=1pt,shorten >=2pt]
            (tmp) -- node[right] {$h$} (S2.center);
    \end{tikzpicture}
\end{figure}
    
\begin{enumerate}
    \item Calculer la valeur $v^{\prime}$ de sa vitesse lorsque son centre s'est
          élevé d'une hauteur $h=\qty{15}{\cm}$.
    \item Quelle distance $L$ parcourt-il le long de la pente avant de s'arrêter
          et de redescendre~?
\end{enumerate}

